\section{信頼性と故障}

\subsection{平成23年度 問13}
\label{q:h23-reg-apt-13}

\noindent\textbf{原典：} sources/H23/aptitude.pdf、PDF物理14頁。

\textbf{公式正答：} 3（single）。

\par\textit{原典本文の校正と独自解説は未作成。分類のみの登録項目。}

\section{公益・説明責任}

\subsection{平成23年度 問15}
\label{q:h23-reg-apt-15}

\noindent\textbf{原典：} sources/H23/aptitude.pdf、PDF物理16頁。

\textbf{公式正答：} 1（single）。

\paragraph{問題の要旨} 社会的合意を形成する際の進行役・専門家・第三者の役割から、不適切な説明を選ぶ。

\Needspace{4\baselineskip}\paragraph{確認した条件・記述（独自要約）}\mbox{}\par

\noindent\textbf{方向：} 最も不適切なもの。\par

\noindent\textbf{状況：} 利害・価値観の異なる関係者間で社会的合意を形成する。\par

\paragraph{問われる核心} 合意形成における専門助言と中立な手続支援

\paragraph{必要知識} 関係者分析、利害、進行役、第三者の中立性、技術的判断と価値判断。

\paragraph{考え方・正答を選ぶ理由（査読待ち初稿）} 合意形成は専門家が一方的に正解を説明する行為ではなく、情報・利害・価値・手続を調整する過程である。担当者の専門性と進行役の中立性・能力を同一視した1が不適切である。第三者を置けば必ず全員一致するという意味でもない。

\Needspace{4\baselineskip}\paragraph{選択肢の検討（査読待ち初稿）}\mbox{}\par

\noindent\textbf{1：設問条件に一致} 業務担当技術者自身が、あらゆる問題に対応する進行役となる。。担当者は事業への利害も持つ。専門知識だけで中立な進行能力やあらゆる対立の処理を保証できない。\par

\noindent\textbf{2：設問条件に不一致} 技術的役割と会議進行・意見調整の役割を分けて検討する。。科学技術の助言と公正な発言機会の確保は異なる能力・役割である。\par

\noindent\textbf{3：設問条件に不一致} 関係者とその利害を把握する。。表面の賛否だけでなく、守りたい利益や懸念を把握して調整する。\par

\noindent\textbf{4：設問条件に不一致} 中立な第三者の支援を得る。。中立性を担保した第三者は手続の公正さを支援し得る。\par

\noindent\textbf{5：設問条件に不一致} 進行支援の必要性を事前評価する。。関係者・対立点・実施可能性を調べて適切な支援を設計する。\par

\paragraph{再利用する考え方} 合意形成における専門助言と中立な手続支援について、対象・要件・例外・時点を分け、類似用語や無条件の断定を避ける。

\paragraph{共通理論} \hyperref[theory:compliance:consensus]{第\ref*{theory:compliance:consensus}節}

\par 参考：\cite{mlit_mediation2006,mlit_pi_thirdparty}。

\par 難易度 3（校準初稿）、優先度未評価。

\subsection{平成28年度 問6}
\label{q:h28-reg-apt-06}

\noindent\textbf{原典：} sources/H28/aptitude.pdf、PDF物理6頁。

\textbf{公式正答：} 1（single）。

\paragraph{問題の要旨} 公衆への説明責任について4記述の正誤の組を選ぶ。

\Needspace{4\baselineskip}\paragraph{確認した条件・記述（独自要約）}\mbox{}\par

\noindent\textbf{ア：} 十分な説明に基づく理解・信頼が重要。\par

\noindent\textbf{イ：} 専門知識を持たない公衆に分かるよう専門家が説明する必要がある。\par

\noindent\textbf{ウ：} 説明責任と守秘義務の双方が必要。\par

\noindent\textbf{エ：} 知る権利に応答する説明を、自分の意見の押付けと区別する。\par

\noindent\textbf{並び：} ア／イ／ウ／エ\par

\paragraph{問われる核心} 専門知識の非対称性と理解可能な説明

\paragraph{必要知識} 説明責任、信頼、知る権利、秘密保持、同意と意見の押付け。

\paragraph{考え方・正答を選ぶ理由（査読待ち初稿）} 4項とも適切であり1である。イは専門外の人にも理解できる説明を行う責任で、同じ専門能力を要求したり、公衆の価値判断を無視したりする趣旨ではない。説明には理解・質問・検討の機会が伴う。守秘を無条件の隠蔽へ、説明を無制限の開示へ変えない。

\Needspace{4\baselineskip}\paragraph{選択肢の検討（査読待ち初稿）}\mbox{}\par

\noindent\textbf{1：設問条件に一致} 正／正／正／正。全項の判断に一致する。\par

\noindent\textbf{2：設問条件に不一致} 誤／誤／正／正。ア・イの正誤が違う。各記述の定義・条件を先に判定してから表に対応させる。\par

\noindent\textbf{3：設問条件に不一致} 正／正／誤／正。ウの正誤が違う。各記述の定義・条件を先に判定してから表に対応させる。\par

\noindent\textbf{4：設問条件に不一致} 正／正／誤／誤。ウ・エの正誤が違う。各記述の定義・条件を先に判定してから表に対応させる。\par

\noindent\textbf{5：設問条件に不一致} 正／誤／誤／正。イ・ウの正誤が違う。各記述の定義・条件を先に判定してから表に対応させる。\par

\paragraph{再利用する考え方} 専門知識の非対称性と理解可能な説明について、対象・要件・例外・時点を分け、類似用語や無条件の断定を避ける。

\paragraph{共通理論} \hyperref[theory:compliance:accountability]{第\ref*{theory:compliance:accountability}節}

\par 参考：\cite{ipej_ethics2011,ieej_conduct2007}。

\par 難易度 3（校準初稿）、優先度未評価。

\section{知的財産}

\subsection{平成27年度 問12}
\label{q:h27-reg-apt-12}

\noindent\textbf{原典：} sources/H27/aptitude.pdf、PDF物理12頁。

\textbf{公式正答：} 2（single）。

\paragraph{問題の要旨} 学術論文等の引用について適切な説明を選ぶ。

\Needspace{4\baselineskip}\paragraph{確認した条件・記述（独自要約）}\mbox{}\par

\noindent\textbf{時点：} 平成27年。\par

\noindent\textbf{方向：} 最も適切なもの。\par

\paragraph{問われる核心} 引用という要件付きの権利制限

\paragraph{必要知識} 著作権法32・43・48条、公表、公正な慣行、正当な範囲、出所、翻訳と翻案。

\paragraph{考え方・正答を選ぶ理由（査読待ち初稿）} 引用は権利者の許諾を必要とする利用一般に対する要件付きの例外であり2が適切である。公表されたから自由という説明、引用だから量・目的を問わないという説明はともに誤る。訳文の利用も引用要件を満たす必要がある。

\Needspace{4\baselineskip}\paragraph{選択肢の検討（査読待ち初稿）}\mbox{}\par

\noindent\textbf{1：設問条件に不一致} 承諾のない引用は全て著作権侵害である。。著作権法32条の要件を満たす引用という権利制限がある。\par

\noindent\textbf{2：設問条件に一致} 適法な引用の条件を満たせば、承諾なしに引用できる。。公表、公正な慣行、目的上正当な範囲などを満たす場合が対象で、出所も明示する。\par

\noindent\textbf{3：設問条件に不一致} 公表済みの著作物は自由に複製できる。。公表と権利放棄は別であり、一般に無制限の複製を許さない。\par

\noindent\textbf{4：設問条件に不一致} 引用は学術研究目的だけに限られる。。報道・批評・研究その他の引用目的で、研究だけとの限定は誤り。\par

\noindent\textbf{5：設問条件に不一致} 外国語の著作物は原語のままでしか引用できない。。適法引用の範囲では43条による翻訳利用が可能であり、原語に限定されない。\par

\paragraph{再利用する考え方} 引用という要件付きの権利制限について、対象・要件・例外・時点を分け、類似用語や無条件の断定を避ける。

\paragraph{共通理論} \hyperref[theory:compliance:citation]{第\ref*{theory:compliance:citation}節}

\par 参考：\cite{bunka_citation_faq,kyushu_copyright2015}。

\paragraph{関連問題} \hyperref[q:r06-reg-apt-03]{R06-REG-APT-03}

\par 難易度 3（校準初稿）、優先度未評価。

\subsection{平成28年度 問11}
\label{q:h28-reg-apt-11}

\noindent\textbf{原典：} sources/H28/aptitude.pdf、PDF物理11頁。

\textbf{公式正答：} 4（single）。

\paragraph{問題の要旨} 知的財産権・特許の発明・実用新案の考案・ソフトウェアの4説明の組を選ぶ。

\Needspace{4\baselineskip}\paragraph{確認した条件・記述（独自要約）}\mbox{}\par

\noindent\textbf{ア：} 法令に定める権利や法律上保護される利益に係る権利を含む知的財産権の定義。\par

\noindent\textbf{イ：} 発明は自然法則を利用した技術的思想の創作のうち高度のもの。\par

\noindent\textbf{ウ：} 考案も自然法則を利用した技術的思想の創作で、高度を要求しない。\par

\noindent\textbf{エ：} 抽象的な人為的取決め・暗号作成・計算は発明に当たらず、プログラムでも自然法則を利用する必要がある。単なるプログラムリストの提示は発明に当たらない。\par

\noindent\textbf{並び：} ア／イ／ウ／エ。平成28年。\par

\paragraph{問われる核心} 技術的創作と単なる規則・情報の区別

\paragraph{必要知識} 知的財産基本法2条、特許法2条、実用新案法2・3条、自然法則、ソフトウェア発明。

\paragraph{考え方・正答を選ぶ理由（査読待ち初稿）} 4項とも基本定義に沿い4である。エの除外は抽象的な規則・数学・情報の単なる提示についてであり、暗号や算法を使う装置・ソフトウェア全般の特許を禁止するものではない。請求項全体としての自然法則の利用を判断する。発明に該当するだけで直ちに特許になるわけではなく、新規性・進歩性等も別要件である。

\Needspace{4\baselineskip}\paragraph{選択肢の検討（査読待ち初稿）}\mbox{}\par

\noindent\textbf{1：設問条件に不一致} 正／正／正／誤。エの正誤が違う。各記述の定義・条件を先に判定してから表に対応させる。\par

\noindent\textbf{2：設問条件に不一致} 正／誤／誤／正。イ・ウの正誤が違う。各記述の定義・条件を先に判定してから表に対応させる。\par

\noindent\textbf{3：設問条件に不一致} 正／誤／誤／誤。イ・ウ・エの正誤が違う。各記述の定義・条件を先に判定してから表に対応させる。\par

\noindent\textbf{4：設問条件に一致} 正／正／正／正。全項の判断に一致する。\par

\noindent\textbf{5：設問条件に不一致} 誤／正／正／誤。ア・エの正誤が違う。各記述の定義・条件を先に判定してから表に対応させる。\par

\paragraph{再利用する考え方} 技術的創作と単なる規則・情報の区別について、対象・要件・例外・時点を分け、類似用語や無条件の断定を避ける。

\paragraph{共通理論} \hyperref[theory:compliance:ip]{第\ref*{theory:compliance:ip}節}, \hyperref[theory:compliance:invention]{第\ref*{theory:compliance:invention}節}

\par 参考：\cite{ip_basic_act,jpo_patent_utility,jpo_invention_criteria}。

\par 難易度 3（校準初稿）、優先度未評価。

\subsection{令和4年度 問9}
\label{q:r04-reg-apt-09}

\noindent\textbf{原典：} sources/R04/aptitude.pdf、PDF物理8頁。

\textbf{公式正答：} 5（single）。

\paragraph{問題の要旨} 知的財産のうち創作物に関する5例の正誤の組を選ぶ。

\Needspace{4\baselineskip}\paragraph{確認した条件・記述（独自要約）}\mbox{}\par

\noindent\textbf{ア：} 発明（特許法）。\par

\noindent\textbf{イ：} 考案（実用新案法）。\par

\noindent\textbf{ウ：} 意匠（意匠法）。\par

\noindent\textbf{エ：} 著作物（著作権法）。\par

\noindent\textbf{オ：} 営業秘密（不正競争防止法）。\par

\noindent\textbf{並び：} ア／イ／ウ／エ／オ。「産業財産権」の数ではなく創作物の区分。\par

\paragraph{問われる核心} 知的創作物と産業財産権の範囲の違い

\paragraph{必要知識} 発明・考案・意匠・著作物・営業秘密、創作物と標識、登録権利と保護利益。

\paragraph{考え方・正答を選ぶ理由（査読待ち初稿）} 5例は全て知的創作物の保護に関わるため5である。営業秘密は秘密管理性・有用性・非公知性を満たす情報で、不正競争防止法が保護する。特許庁に登録して生じる4種類の産業財産権に限定すると、著作物や営業秘密を誤って除外してしまう。

\Needspace{4\baselineskip}\paragraph{選択肢の検討（査読待ち初稿）}\mbox{}\par

\noindent\textbf{1：設問条件に不一致} 正／誤／正／正／正。イの正誤が違う。各記述の定義・条件を先に判定してから表に対応させる。\par

\noindent\textbf{2：設問条件に不一致} 正／正／誤／正／正。ウの正誤が違う。各記述の定義・条件を先に判定してから表に対応させる。\par

\noindent\textbf{3：設問条件に不一致} 正／正／正／誤／正。エの正誤が違う。各記述の定義・条件を先に判定してから表に対応させる。\par

\noindent\textbf{4：設問条件に不一致} 正／正／正／正／誤。オの正誤が違う。各記述の定義・条件を先に判定してから表に対応させる。\par

\noindent\textbf{5：設問条件に一致} 正／正／正／正／正。全項の判断に一致する。\par

\paragraph{再利用する考え方} 知的創作物と産業財産権の範囲の違いについて、対象・要件・例外・時点を分け、類似用語や無条件の断定を避ける。

\paragraph{共通理論} \hyperref[theory:compliance:ip]{第\ref*{theory:compliance:ip}節}, \hyperref[theory:compliance:secret]{第\ref*{theory:compliance:secret}節}

\par 参考：\cite{ip_basic_act,jpo_ip_overview,meti_trade_secret}。

\paragraph{関連問題} \hyperref[q:r05-reg-apt-04]{R05-REG-APT-04}

\par 難易度 3（校準初稿）、優先度未評価。

\subsection{令和5年度 問4}
\label{q:r05-reg-apt-04}

\noindent\textbf{原典：} sources/R05/aptitude.pdf、PDF物理4頁。

\textbf{公式正答：} 2（single）。

\paragraph{問題の要旨} 産業財産権に含まれる5種類の権利の正誤の組を選ぶ。

\Needspace{4\baselineskip}\paragraph{確認した条件・記述（独自要約）}\mbox{}\par

\noindent\textbf{並び：} ア：特許権／イ：実用新案権／ウ：意匠権／エ：商標権／オ：著作権。\par

\paragraph{問われる核心} 産業財産権4種類と広い知的財産権の区別

\paragraph{必要知識} 特許庁所管4権利、著作権、権利の対象と制度の範囲。

\paragraph{考え方・正答を選ぶ理由（査読待ち初稿）} 産業財産権は特許・実用新案・意匠・商標の4種類である。著作権も知的財産権であるが、この4種類に含まれない。正／正／正／正／誤の2を選ぶ。産業に使われる著作物だから産業財産権になるという分類ではない。

\Needspace{4\baselineskip}\paragraph{選択肢の検討（査読待ち初稿）}\mbox{}\par

\noindent\textbf{1：設問条件に不一致} 正／正／正／正／正。オの正誤が違う。各記述の定義・条件を先に判定してから表に対応させる。\par

\noindent\textbf{2：設問条件に一致} 正／正／正／正／誤。全項の判断に一致する。\par

\noindent\textbf{3：設問条件に不一致} 正／正／正／誤／正。エ・オの正誤が違う。各記述の定義・条件を先に判定してから表に対応させる。\par

\noindent\textbf{4：設問条件に不一致} 正／正／誤／正／正。ウ・オの正誤が違う。各記述の定義・条件を先に判定してから表に対応させる。\par

\noindent\textbf{5：設問条件に不一致} 正／誤／正／正／正。イ・オの正誤が違う。各記述の定義・条件を先に判定してから表に対応させる。\par

\paragraph{再利用する考え方} 産業財産権4種類と広い知的財産権の区別について、対象・要件・例外・時点を分け、類似用語や無条件の断定を避ける。

\paragraph{共通理論} \hyperref[theory:compliance:ip]{第\ref*{theory:compliance:ip}節}

\par 参考：\cite{jpo_ip_overview}。

\paragraph{関連問題} \hyperref[q:r04-reg-apt-09]{R04-REG-APT-09}

\par 難易度 3（校準初稿）、優先度未評価。

\subsection{令和6年度 問3}
\label{q:r06-reg-apt-03}

\noindent\textbf{原典：} sources/R06/aptitude.pdf、PDF物理3頁。

\textbf{公式正答：} 1（single）。

\paragraph{問題の要旨} 引用・官公資料・AI学習・技術試験の著作権制限について正誤の組を選ぶ。

\Needspace{4\baselineskip}\paragraph{確認した条件・記述（独自要約）}\mbox{}\par

\noindent\textbf{ア：} 学術論文を無承諾で引用すれば全て侵害となる。\par

\noindent\textbf{イ：} 広く利用すべき官公資料は、転載禁止表示がない等の要件の下で説明材料に転載できる。\par

\noindent\textbf{ウ：} AI開発目的の学習データ収集や開発する第三者への提供は権利制限の対象となる。\par

\noindent\textbf{エ：} 複製向けのカメラ・プリンター・和紙等を開発するための試験的複製は権利制限の対象となる。\par

\noindent\textbf{並び：} ア／イ／ウ／エ。2024年。必要な限度と権利者の利益を不当に害しない条件を適用。\par

\paragraph{問われる核心} 著作権保護と利用を調整する要件付きの権利制限

\paragraph{必要知識} 著作権法32条・30条の4、引用、官公資料、非享受目的、必要限度、不当な利益侵害、2019年施行改正。

\paragraph{考え方・正答を選ぶ理由（査読待ち初稿）} アは32条の適法引用を無視しているので誤り。イは32条2項の資料・刊行物等の要件内、ウ・エは30条の4の非享受目的・必要な限度・不当な利益侵害をしない範囲の例で正しい。誤／正／正／正の1である。AI学習に使うと言えば常に自由になるわけではなく、作品の表現を出力させる享受目的が併存する利用や解析用データベースの市場を害する利用は別に検討する。生成した出力の利用も学習段階とは別である。

\Needspace{4\baselineskip}\paragraph{選択肢の検討（査読待ち初稿）}\mbox{}\par

\noindent\textbf{1：設問条件に一致} 誤／正／正／正。全項の判断に一致する。\par

\noindent\textbf{2：設問条件に不一致} 正／誤／正／正。ア・イの正誤が違う。各記述の定義・条件を先に判定してから表に対応させる。\par

\noindent\textbf{3：設問条件に不一致} 正／正／誤／正。ア・ウの正誤が違う。各記述の定義・条件を先に判定してから表に対応させる。\par

\noindent\textbf{4：設問条件に不一致} 正／正／正／誤。ア・エの正誤が違う。各記述の定義・条件を先に判定してから表に対応させる。\par

\noindent\textbf{5：設問条件に不一致} 正／正／正／正。アの正誤が違う。各記述の定義・条件を先に判定してから表に対応させる。\par

\paragraph{再利用する考え方} 著作権保護と利用を調整する要件付きの権利制限について、対象・要件・例外・時点を分け、類似用語や無条件の断定を避ける。

\paragraph{共通理論} \hyperref[theory:compliance:citation]{第\ref*{theory:compliance:citation}節}, \hyperref[theory:compliance:nonenjoyment]{第\ref*{theory:compliance:nonenjoyment}節}

\par 参考：\cite{bunka_citation_faq,bunka_amend2018,bunka_ai2024,bunka_text2026}。

\paragraph{関連問題} \hyperref[q:h27-reg-apt-12]{H27-REG-APT-12}

\par 難易度 3（校準初稿）、優先度未評価。

\section{公正取引・競争法・金融市場}

\subsection{平成27年度 問10}
\label{q:h27-reg-apt-10}

\noindent\textbf{原典：} sources/H27/aptitude.pdf、PDF物理10頁。

\textbf{公式正答：} 5（single）。

\paragraph{問題の要旨} 公共・民間入札の契約方式と談合への法的対応から、不適切な説明を選ぶ。

\Needspace{4\baselineskip}\paragraph{確認した条件・記述（独自要約）}\mbox{}\par

\noindent\textbf{時点：} 平成27年の制度。\par

\noindent\textbf{方向：} 最も不適切なもの。\par

\paragraph{問われる核心} 公共契約方式と発注主体に依存しない競争法規制

\paragraph{必要知識} 一般競争・随意契約、独占禁止法、不当な取引制限、行政・刑事責任、官製談合。

\paragraph{考え方・正答を選ぶ理由（査読待ち初稿）} 発注者が公的か民間かという区別と、事業者間で競争を制限したかという区別は別である。民間入札なら適法という5の断定が誤りである。公共契約の例外、競争法上の違反、職員の関与の責任も同じ制度へまとめない。

\Needspace{4\baselineskip}\paragraph{選択肢の検討（査読待ち初稿）}\mbox{}\par

\noindent\textbf{1：設問条件に不一致} 公共調達は一般競争入札が基本である。。国の会計法上、一般競争を原則とする。地方公共団体等の具体的根拠は各制度に従う。\par

\noindent\textbf{2：設問条件に不一致} 法定条件では随意契約が認められる。。性質・目的が競争を許さない、緊急等の法定要件があり、自由に相手を選ぶ例外ではない。\par

\noindent\textbf{3：設問条件に不一致} 入札談合に行政処分と刑事罰が科され得る。。独占禁止法上の措置や課徴金、刑事処罰の対象となり得る。ただし各行為・責任の要件があり、自動的に全参加者へ同じ制裁が及ぶとはしない。\par

\noindent\textbf{4：設問条件に不一致} 官製談合への職員の関与に刑事罰の制度がある。。官製談合防止制度には一定の職員の入札等妨害への刑事罰があり、2015年以前に施行されている。\par

\noindent\textbf{5：設問条件に一致} 民間事業者による入札の談合は独占禁止法等に抵触しない。。独占禁止法は民間調達の競争制限を一律に除外しない。民間の受注調整にも実際の法的措置がある。\par

\paragraph{再利用する考え方} 公共契約方式と発注主体に依存しない競争法規制について、対象・要件・例外・時点を分け、類似用語や無条件の断定を避ける。

\paragraph{共通理論} \hyperref[theory:compliance:procurement]{第\ref*{theory:compliance:procurement}節}, \hyperref[theory:compliance:competition]{第\ref*{theory:compliance:competition}節}

\par 参考：\cite{mof_procurement2006,jftc_rules,jftc_kansei,jftc_private2018}。

\par 難易度 3（校準初稿）、優先度未評価。

\subsection{令和4年度 問12}
\label{q:r04-reg-apt-12}

\noindent\textbf{原典：} sources/R04/aptitude.pdf、PDF物理11頁。

\textbf{公式正答：} 5（single）。

\paragraph{問題の要旨} 公正取引を損なう4行為と名称について正誤の組を選ぶ。

\Needspace{4\baselineskip}\paragraph{確認した条件・記述（独自要約）}\mbox{}\par

\noindent\textbf{ア：} ア：公共入札で事前に受注者等を決める行為をインサイダー取引という。\par

\noindent\textbf{イ：} イ：自然な需給に見せて人為的に価格を動かし利益を図る行為を相場操縦という。\par

\noindent\textbf{ウ：} ウ：事業者間で商品価格や生産・販売数量を共同決定する行為を談合という。\par

\noindent\textbf{エ：} エ：会社関係者等が未公表の重要事実に関係する株式等の取引をする行為をカルテルという。\par

\paragraph{問われる核心} 競争制限と金融市場の公正を損なう行為の識別

\paragraph{必要知識} カルテル、入札談合、インサイダー取引、相場操縦、規制する行為と法の違い。

\paragraph{考え方・正答を選ぶ理由（査読待ち初稿）} 入札の受注者・価格の事前調整は入札談合、商品の価格・数量の共同決定はカルテル、未公表の重要事実に基づく関係者等の規制はインサイダー取引、人為的な価格形成で他者を誘引する行為は相場操縦である。誤／正／誤／誤の5。なお、入札談合は競争法上カルテルの典型例でもある。用語の出題上の区別を、法的な概念が全く重ならないという主張にしない。

\Needspace{4\baselineskip}\paragraph{選択肢の検討（査読待ち初稿）}\mbox{}\par

\noindent\textbf{1：設問条件に不一致} 正／誤／正／正。ア・イ・ウ・エの正誤が違う。各記述の定義・条件を先に判定してから表に対応させる。\par

\noindent\textbf{2：設問条件に不一致} 正／正／正／誤。ア・ウの正誤が違う。各記述の定義・条件を先に判定してから表に対応させる。\par

\noindent\textbf{3：設問条件に不一致} 誤／正／誤／正。エの正誤が違う。各記述の定義・条件を先に判定してから表に対応させる。\par

\noindent\textbf{4：設問条件に不一致} 正／誤／誤／正。ア・イ・エの正誤が違う。各記述の定義・条件を先に判定してから表に対応させる。\par

\noindent\textbf{5：設問条件に一致} 誤／正／誤／誤。全項の判断に一致する。\par

\paragraph{再利用する考え方} 競争制限と金融市場の公正を損なう行為の識別について、対象・要件・例外・時点を分け、類似用語や無条件の断定を避ける。

\paragraph{共通理論} \hyperref[theory:compliance:competition]{第\ref*{theory:compliance:competition}節}, \hyperref[theory:compliance:market]{第\ref*{theory:compliance:market}節}

\par 参考：\cite{jftc_rules,jftc_bid_cartel2021,fsa_market2004,fsa_information}。

\paragraph{関連問題} \hyperref[q:r07-reg-apt-11]{R07-REG-APT-11}

\par 難易度 3（校準初稿）、優先度未評価。

\subsection{令和7年度 問11}
\label{q:r07-reg-apt-11}

\noindent\textbf{原典：} sources/R07/aptitude.pdf、PDF物理11頁。

\textbf{公式正答：} 4（single）。

\paragraph{問題の要旨} 公正取引を損なう4行為と名称について適切な記述の数を選ぶ。

\Needspace{4\baselineskip}\paragraph{確認した条件・記述（独自要約）}\mbox{}\par

\noindent\textbf{ア：} ア：事業者間で商品価格や生産・販売数量を共同決定する行為を談合という。\par

\noindent\textbf{イ：} イ：会社関係者等が未公表の重要事実に関係する株式等の取引をする行為をカルテルという。\par

\noindent\textbf{ウ：} ウ：自然な需給に見せて人為的に価格を動かし利益を図る行為を相場操縦という。\par

\noindent\textbf{エ：} エ：公共入札で事前に受注者等を決める行為をインサイダー取引という。\par

\paragraph{問われる核心} 競争制限と金融市場の公正を損なう行為の識別

\paragraph{必要知識} カルテル、入札談合、インサイダー取引、相場操縦、規制する行為と法の違い。

\paragraph{考え方・正答を選ぶ理由（査読待ち初稿）} 入札の受注者・価格の事前調整は入札談合、商品の価格・数量の共同決定はカルテル、未公表の重要事実に基づく関係者等の規制はインサイダー取引、人為的な価格形成で他者を誘引する行為は相場操縦である。ウだけ正しく1項なので4。なお、入札談合は競争法上カルテルの典型例でもある。用語の出題上の区別を、法的な概念が全く重ならないという主張にしない。

\Needspace{4\baselineskip}\paragraph{選択肢の検討（査読待ち初稿）}\mbox{}\par

\noindent\textbf{1：設問条件に不一致} 4。適切なのはウだけで、4項ではない。\par

\noindent\textbf{2：設問条件に不一致} 3。アはカルテル、イはインサイダー、エは入札談合なので3項ではない。\par

\noindent\textbf{3：設問条件に不一致} 2。ウ以外は行為と名称が合わず、2項ではない。\par

\noindent\textbf{4：設問条件に一致} 1。相場操縦のウだけが適切で、1項である。\par

\noindent\textbf{5：設問条件に不一致} 0。ウの説明が適切なので0項ではない。\par

\paragraph{再利用する考え方} 競争制限と金融市場の公正を損なう行為の識別について、対象・要件・例外・時点を分け、類似用語や無条件の断定を避ける。

\paragraph{共通理論} \hyperref[theory:compliance:competition]{第\ref*{theory:compliance:competition}節}, \hyperref[theory:compliance:market]{第\ref*{theory:compliance:market}節}

\par 参考：\cite{jftc_rules,jftc_bid_cartel2021,fsa_market2004,fsa_information}。

\paragraph{関連問題} \hyperref[q:r04-reg-apt-12]{R04-REG-APT-12}

\par 難易度 3（校準初稿）、優先度未評価。

\section{社会的責任・安全の標準}

\subsection{令和6年度 問12}
\label{q:r06-reg-apt-12}

\noindent\textbf{原典：} sources/R06/aptitude.pdf、PDF物理12頁。

\textbf{公式正答：} 3（single）。

\paragraph{問題の要旨} ISO・IEC・JISの名称、認証、審査、図の新旧の6説明について適切なものの数を選ぶ。

\Needspace{4\baselineskip}\paragraph{確認した条件・記述（独自要約）}\mbox{}\par

\noindent\textbf{ア：} ISOは電気・通信・電子技術分野を除く産業分野の国際規格を作る。\par

\noindent\textbf{イ：} IECは電気・電子技術の国際規格を作る。\par

\noindent\textbf{ウ：} 2019年に産業標準化法→工業標準化法、日本産業規格→日本工業規格と改称した。\par

\noindent\textbf{エ：} 国登録の民間第三者機関の適合性審査と認証でJISマークを表示できる。\par

\noindent\textbf{オ：} 登録認証機関は少なくとも3年ごと・必要時臨時に審査し、国も必要に応じ立入検査する。\par

\noindent\textbf{カ：} 図の左を旧、右を現行JISマークと説明している。\par

\paragraph{問われる核心} 国際規格の分担・JISの改称・認証と図の新旧

\paragraph{必要知識} ISO、IEC、産業標準化法、2019年改称、登録認証機関、定期審査、2005年新マーク。

\paragraph{考え方・正答を選ぶ理由（査読待ち初稿）} JISCの分担説明からア・イは正しい。2019年7月1日の改称は工業→産業なのでウは誤り。エ・オは認証・審査制度に沿う。図では左の横並びJISが2005年開始の現行、右が旧で、カの説明は逆である。4項が正しく3を選ぶ。初回独立解答2からの修正は公式番号だけで行わず、JISCの現行マーク画像を別に表示して照合し、その画像・ハッシュ・変更理由と初回記録を保存している。2019年の対象拡大と2005年のマーク制度変更を混同しない。

\Needspace{4\baselineskip}\paragraph{選択肢の検討（査読待ち初稿）}\mbox{}\par

\noindent\textbf{1：設問条件に不一致} 6。ウは改称方向が逆、カは図の新旧が逆なので6項ではない。\par

\noindent\textbf{2：設問条件に不一致} 5。初回の独立判断ではカを正として5項としたが、JISC公式現行図と実画像を照合しカも誤りと確定した。\par

\noindent\textbf{3：設問条件に一致} 4。ア・イ・エ・オの4項が適切である。\par

\noindent\textbf{4：設問条件に不一致} 3。名称と図以外の4項は適切なので3項ではない。\par

\noindent\textbf{5：設問条件に不一致} 2。ウとカが誤りであり、適切なのは2項ではなく4項である。\par

\paragraph{再利用する考え方} 国際規格の分担・JISの改称・認証と図の新旧について、対象・要件・例外・時点を分け、類似用語や無条件の断定を避ける。

\paragraph{共通理論} \hyperref[theory:compliance:standards]{第\ref*{theory:compliance:standards}節}

\par 参考：\cite{jisc_isoiec,jisc_certification,jisc_marks,meti_jis2019,mhlw_jis_history}。

\par 難易度 3（校準初稿）、優先度未評価。

\section{公益通報者保護と通報経路}

\subsection{平成30年度 問8}
\label{q:h30-reg-apt-08}

\noindent\textbf{原典：} sources/H30/aptitude.pdf、PDF物理8頁。

\textbf{公式正答：} 3（single）。

\paragraph{問題の要旨} 出題当時の公益通報者保護法について、不適切な説明を選ぶ。

\Needspace{4\baselineskip}\paragraph{確認した条件・記述（独自要約）}\mbox{}\par

\noindent\textbf{時点：} 2018年。2020年改正・2022年施行前。\par

\noindent\textbf{方向：} 最も不適切なもの。\par

\paragraph{問われる核心} 公益通報の対象者・通報先別の保護要件と施行時点

\paragraph{必要知識} 公益通報者保護法2・3・7条、不正目的、相当の理由、公務員、内部・行政・外部、追加条件。

\paragraph{考え方・正答を選ぶ理由（査読待ち初稿）} 公務員を一律に除外する3が誤りである。通報先によって保護要件が異なり、その他外部への通報には当時の3条3号の条件を確認する。原典5の「など」は例示なので、内部通報後20日経過等と急迫した生命・身体の危険を全て同時に要求するとは解さない。現在は2022年施行改正で保護対象・体制・要件が更新されている。2025年改正の主要部分は2026年12月1日施行予定であり、2026年10月10日現在の施行法と混ぜない。

\Needspace{4\baselineskip}\paragraph{選択肢の検討（査読待ち初稿）}\mbox{}\par

\noindent\textbf{1：設問条件に不一致} 不正目的でなく通報対象事実を所定の通報先へ知らせる。。不正の目的を排除し、法律が定める対象事実と通報先を満たす必要がある。\par

\noindent\textbf{2：設問条件に不一致} 保護要件を満たす通報者の解雇・不利益取扱いを防ぐ。。通報を理由とする不利益から保護する制度で、あらゆる行為を無条件に免責するものではない。\par

\noindent\textbf{3：設問条件に一致} 保護対象の労働者に公務員は含まれない。。公務員も原則として労働者に含まれるので誤り。公務員への取扱いには法の特則がある。\par

\noindent\textbf{4：設問条件に不一致} 内部・行政・その他外部で保護要件が異なる。。当時、行政への通報では相当の理由、その他外部ではさらに所定の追加条件などが必要であった。\par

\noindent\textbf{5：設問条件に不一致} その他外部への通報には相当の理由と被害防止に必要な通報先、内部調査通知がない等の追加条件が求められる。。列挙された事情は法定の選択的条件の例示である。全条件の同時成立や、必ず内部を経由する義務とは読まない。\par

\paragraph{再利用する考え方} 公益通報の対象者・通報先別の保護要件と施行時点について、対象・要件・例外・時点を分け、類似用語や無条件の断定を避ける。

\paragraph{共通理論} \hyperref[theory:compliance:whistleblower]{第\ref*{theory:compliance:whistleblower}節}

\par 参考：\cite{caa_handbook_pre2022,caa_report_current,caa_amend2020,caa_schedule2026}。

\par 難易度 3（校準初稿）、優先度未評価。
